\documentclass[10pt,a4paper,twocolumn]{article}

\usepackage[hyphens]{url}
\usepackage[margin=1.8cm]{geometry}
\usepackage{amsmath,amssymb}
\usepackage{booktabs}
\usepackage{array}
\usepackage{hyperref}
\usepackage{xcolor}
\usepackage{caption}
\usepackage{enumitem}
\usepackage{microtype}
\usepackage{float}
\usepackage{titlesec}
\usepackage{balance}

\titleformat{\section}{\normalsize\bfseries}{\thesection.}{0.5em}{}
\titleformat{\subsection}{\small\bfseries}{\thesubsection}{0.5em}{}
\setlength{\columnsep}{0.6cm}

\hypersetup{colorlinks=true,linkcolor=blue!60!black,citecolor=blue!60!black,urlcolor=blue!50!black}

\begin{document}

% ── Title ─────────────────────────────────────────────────────────────────────
\twocolumn[{%
\centering
{\large\bfseries Methodology: Stage-Aware IT Power Modeling\\
and Validation on Marconi100\par}
\vspace{0.3em}
{\small Prince Tawiah \quad Michigan Technological University}

\vspace{0.8em}
\begin{minipage}{0.92\linewidth}
\small\textbf{Abstract.}
This document describes the methodology implemented in the IT-power simulation
and validation toolchain. Fleet GPU power is computed from a Fan-style linear
device model extended with an idle fraction~$\rho_k$ and the full installed
inventory~$N_k$. Per-job utilization $u_m(t)$ is stage-dependent (training,
fine-tuning, inference), with waveforms motivated by Li et~al.\ and related
AI-power studies. On Marconi100 (M100) ExaData, SLURM \texttt{job\_table}
schedules are replayed, stages are assigned by duration--scale heuristics, and
modeled fleet GPU power is compared to Ganglia \texttt{Gpu*\_power\_usage}.
An interactive Schedule Explorer (\texttt{validation\_v2}) and a synthetic-task
simulator (\texttt{server.py}) operationalize the method.
\end{minipage}
\vspace{0.8em}
}]

% ── 1. Scope ──────────────────────────────────────────────────────────────────
\section{Scope and Objectives}

The methodology supports two complementary goals:

\begin{enumerate}[leftmargin=1.2em,itemsep=2pt,topsep=2pt]
  \item \textbf{Forward simulation:} given a set of AI tasks (manual or
        uploaded), estimate time-varying IT/GPU power by hardware type and
        lifecycle stage.
  \item \textbf{Replay validation:} given real M100 job schedules and Ganglia
        GPU telemetry for a fixed window, quantify how closely the same physics
        matches measured fleet GPU power.
\end{enumerate}

Facility-level \texttt{Tot\_ict} (logics) is available in ExaData but is
\emph{not} the primary target of the current Schedule Explorer, which validates
\textbf{GPU power only}.

% ── 2. Core power model ───────────────────────────────────────────────────────
\section{Core Power Formula}
\label{sec:formula}

For each hardware class~$k$ (e.g.\ V100), let $N_k$ be the installed device
count, $P_k^{\max}$ the device TDP, $\rho_k$ the idle fraction of TDP, and
$\eta_k(t)$ the fraction of the inventory allocated to active jobs. The fleet
macro power (kW when $P$ is in watts) is
\begin{equation}
  P_k(t) = N_k\,P_k^{\max}
  \bigl[\eta_k(t)\,U_k(t) + \bigl(1-\eta_k(t)\bigr)\rho_k\bigr],
  \label{eq:core}
\end{equation}
with total IT contribution
\begin{equation}
  P_{\mathrm{macro}}(t)=\sum_k P_k(t),\qquad
  P_{\mathrm{total}}(t)=P_{\mathrm{macro}}(t)+P_{\mathrm{infra}}.
  \label{eq:total}
\end{equation}
Infrastructure $P_{\mathrm{infra}}$ (network/storage/misc) is an optional
constant offset in the interactive simulator and is set to zero in M100 GPU
validation.

\subsection{Allocation and weighted utilization}

At time~$t$, let $\mathcal{A}_k(t)$ be the set of active tasks on hardware~$k$.
Demand $D_k(t)=\sum_{m\in\mathcal{A}_k} n_m$ (requested devices). If
$D_k\le N_k$, each task keeps its request; otherwise requests are scaled
pro~rata:
\begin{equation}
  n_m^{\mathrm{eff}} = n_m\cdot\min\!\Bigl(1,\frac{N_k}{D_k}\Bigr),\qquad
  N_k^{\mathrm{alloc}}=\min(D_k,N_k),\qquad
  \eta_k=\frac{N_k^{\mathrm{alloc}}}{N_k}.
\end{equation}
The busy utilization is a device-weighted average of per-task utilizations:
\begin{equation}
  U_k(t)=\sum_{m\in\mathcal{A}_k}
  \alpha_m\,u_m(t),\qquad
  \alpha_m=\frac{n_m^{\mathrm{eff}}}{N_k^{\mathrm{alloc}}},
  \label{eq:Uk}
\end{equation}
then floored at idle, $U_k\leftarrow\max(U_k,\rho_k)$. Unallocated devices
contribute only the idle term $(1-\eta_k)\rho_k$ inside~\eqref{eq:core}; no
separate inventory-idle addend is required.

\subsection{Idle fraction $\rho_k$}

We define
\begin{equation}
  \rho_k=\frac{P_{k,\mathrm{idle}}}{P_k^{\max}},
  \label{eq:rho}
\end{equation}
so idle power per device is $\rho_k P_k^{\max}$. This differs from a
capacity-normalized estimator; $\rho_k$ is treated as a hardware constant
in \texttt{config.json}.

\begin{table}[H]
\centering
\caption{Hardware library used by the simulator (from \texttt{config.json}).}
\label{tab:hw}
\renewcommand{\arraystretch}{1.15}
\small
\begin{tabular}{lccc}
\toprule
HW & $P^{\max}$ (W) & $\rho_k$ & Role \\
\midrule
H100    & 700 & 0.129 & library \\
A100    & 400 & 0.288 & library \\
TPUv4   & 450 & 0.176 & library \\
CPU STD & 400 & 0.450 & library \\
CPU HPC & 550 & 0.451 & library \\
V100    & 300 & $\approx$0.117 & \textbf{M100 replay} \\
\bottomrule
\end{tabular}
\end{table}

\noindent For M100 validation the active class is V100 with
$P^{\max}=300\,\mathrm{W}$ (Volta TDP) and four GPUs per node. Inventory for
replay is set near the Ganglia reporting footprint
($N_{\mathrm{V100}}=3780$ in the current preset; cluster nameplate is
980 nodes~$\times$4).

% ── 3. Stages ─────────────────────────────────────────────────────────────────
\section{Stage Utilization Models}
\label{sec:stages}

Jobs/tasks are labeled as \texttt{training}, \texttt{fine\_tuning}, or
\texttt{inference}. Utilization $u_m(t)$ is drawn from stage-specific
stochastic processes implemented in \texttt{stochastic\_util}
(\texttt{server.py}), with parameters in \texttt{stage\_physics} and clamp
bounds in \texttt{stage\_defaults}.

\subsection{Stage assignment on M100 replay}

ExaData \texttt{job\_table} does not provide AI-stage labels. For V100
replay we assign stages by duration and node count
(\texttt{replay.stage\_thresholds}):
\begin{equation}
\begin{aligned}
&\text{training if } T\ge 1\,\mathrm{h}\ \wedge\ n_{\mathrm{nodes}}\ge 16,\\
&\text{inference if } T<0.25\,\mathrm{h}\ \vee\ n_{\mathrm{nodes}}<4,\\
&\text{fine-tuning otherwise.}
\end{aligned}
\end{equation}
This is a \emph{heuristic proxy}, not ground truth; mislabeling is a known
limitation when comparing shapes to Ganglia.

\subsection{Training}

Most of the time the job sits on a compute plateau:
\begin{equation}
  u\sim\mathcal{N}(u_{\mathrm{plateau}},\sigma_{\mathrm{plateau}}^2),
  \quad u\leftarrow\mathrm{clip}(u;u_{\min},u_{\max}).
\end{equation}
Every $\texttt{checkpoint\_ms}$ (default 300\,s) a dip of duration
$\texttt{checkpoint\_dur\_ms}$ draws from a low checkpoint utilization
band. Parameters are motivated by checkpoint/communication structure in
large-model training literature \cite{li2024unseen,shoeybi2019megatron}.

\subsection{Fine-tuning}

Analogous plateau with periodic \emph{evaluation} dips
(\texttt{eval\_period\_ms}, \texttt{eval\_dur\_ms}, low $u_{\mathrm{eval}}$),
following the qualitative shape in Li et~al.\ Fig.~13 \cite{li2024unseen}.

\subsection{Inference}

Without an explicit request schedule (M100 replay path), occupancy is
approximated from rate $\lambda$ (req/h) and service time $\tau$:
\begin{equation}
  p_{\mathrm{busy}}=\min\bigl(1,\lambda\tau/3600\bigr).
\end{equation}
With probability $p_{\mathrm{busy}}$ draw burst utilization
$\mathcal{N}(u_{\mathrm{burst}},\sigma_{\mathrm{burst}}^2)$; otherwise draw
idle utilization floored by $\max(u_{\mathrm{idle}},\rho_k)$ and capped by
$\texttt{idle\_u\_max}$. When a synthetic schedule is available, busy/idle
device counts are taken from Poisson-generated request intervals instead.

\begin{table}[H]
\centering
\caption{Representative stage parameters (tunable in \texttt{config.json}).}
\label{tab:stages}
\small
\renewcommand{\arraystretch}{1.15}
\begin{tabular}{lccc}
\toprule
Stage & Plateau / burst & Clamp / notes & Native micro step \\
\midrule
Training & $u_{\mathrm{plateau}}$ & $u_{\min}$--$u_{\max}$; 5\,min ckpt & 10\,s \\
Fine-tuning & $u_{\mathrm{plateau}}$ & eval dips $\sim$4\,min & 1\,s \\
Inference & $u_{\mathrm{burst}}$ / $u_{\mathrm{idle}}$ & $\lambda$, $\tau$ & 100\,ms \\
\bottomrule
\end{tabular}
\end{table}

\noindent Macro charts in the interactive simulator downsample to at most
\texttt{macro\_max\_pts}; per-stage micro panels illustrate waveform
\emph{shape} and are \textbf{not} required to sum to the fleet macro curve.

% ── 4. Validation on M100 ─────────────────────────────────────────────────────
\section{Validation Methodology (M100 ExaData)}
\label{sec:validation}

\subsection{Data sources}

\begin{itemize}[leftmargin=1.2em,itemsep=1pt,topsep=2pt]
  \item \texttt{plugin=job\_table}: job start/end, \texttt{num\_nodes}, ids
        (metric \texttt{job\_info\_marconi100}).
  \item \texttt{plugin=ganglia\_pub}: per-GPU
        \texttt{Gpu\{0--3\}\_power\_usage} (W).
\end{itemize}
Measured fleet GPU power is formed by flooring samples to 1\,min, averaging
per (minute, node, Gpu metric), summing watts across reporting GPUs, and
converting to kW. This includes idle and busy GPUs that report---there is no
Slurm allocation filter on the measurement.

\subsection{Window and caching}

Default validation window: 2022-03-15\,UTC (configurable in
\texttt{validation\_v2/slice\_config.json}). Reporting grid
$\Delta t_{\mathrm{grid}}=300\,\mathrm{s}$. Modeled power is evaluated on a
finer sub-step $\Delta t_{\mathrm{sub}}$ (default 30\,s) inside each bin and
averaged, so checkpoint/eval dips are not systematically aliased to the
5-minute grid. Results are cached in
\texttt{validation\_v2/slice\_cache/day\_slice.json}; changing
\texttt{config.json} requires rebuilding the slice.

\subsection{Modeled series}

All jobs overlapping the window are converted to tasks
(\texttt{job\_to\_task}), classified, and simulated with~\eqref{eq:core}.
Outputs include: fleet $P(t)$; per-stage attributed power (allocated devices
only); active job/node counts; a filtered Gantt for visualization. Stage
charts need not sum exactly to fleet power because unallocated idle
$(1-\eta)\rho N P^{\max}$ appears only in the fleet curve.

\subsection{Comparison metrics}

Primary comparisons against Ganglia:
mean/peak power, pointwise residual, MAPE, and Pearson correlation~$r$ on
the aligned grid. Level agreement and shape agreement are reported separately:
mean-matched runs can still be anti-correlated if stage/schedule attribution
is wrong.

% ── 5. Software architecture ──────────────────────────────────────────────────
\section{Software Architecture}

\begin{itemize}[leftmargin=1.2em,itemsep=1pt,topsep=2pt]
  \item \texttt{config.json} --- hardware, stage physics, M100 preset, replay
        thresholds.
  \item \texttt{server.py} + \texttt{index.html} --- interactive simulator
        (port 5000): manual tasks, macro/micro charts, optional job replay API.
  \item \texttt{validation/exadata.py} --- Parquet loaders for ExaData.
  \item \texttt{validation\_v2/} --- Schedule Explorer (port 5051):
        measured vs modeled GPU power, Gantt, stage panels.
\end{itemize}

Physics is shared: validation imports \texttt{stochastic\_util},
\texttt{job\_to\_task}, and hardware constants from \texttt{server.py}.

% ── 6. Limitations ────────────────────────────────────────────────────────────
\section{Limitations and Tuning Practice}

\begin{enumerate}[leftmargin=1.2em,itemsep=2pt,topsep=2pt]
  \item Stage labels are heuristic; Ganglia has no stage ground truth.
  \item Inference occupancy without per-job request traces is a mean-field
        approximation ($\lambda,\tau$) and can under-represent sharp spikes.
  \item Clamp bounds $u_{\min},u_{\max}$ must remain consistent with
        $u_{\mathrm{plateau}}$; otherwise plateau edits have no effect.
  \item $P^{\max}$ should remain device TDP (300\,W for V100); substituting
        observed average watts collapses physical meaning of $\rho$ and $u$.
  \item Dynamic terms $C\dot P+D\ddot P$ from earlier drafts are
        \emph{not} in the current validation path; they remain future work.
  \item Practical tuning separates \textbf{level} (util, $\lambda$, $\rho$,
        $N_k$) from \textbf{shape} (stage rules, concurrency response,
        schedule fidelity). Matching day-mean alone is insufficient when~$r$
        is negative.
\end{enumerate}

\balance
\begin{thebibliography}{9}

\bibitem{fan2007power}
W.~Fan \textit{et al.},
``Power provisioning for a warehouse-sized computer,''
\textit{ISCA}, 2007.

\bibitem{li2024unseen}
Y.~Li \textit{et al.},
``The unseen AI disruptions for power grids: LLM-induced transients,''
\textit{arXiv:2409.11416}, 2024.

\bibitem{chen2025electricity}
X.~Chen \textit{et al.},
``Electricity demand and grid impacts of AI data centers,''
\textit{arXiv:2509.07218}, 2025.

\bibitem{choukse2025power}
E.~Choukse \textit{et al.},
``Power stabilization for AI training datacenters,''
\textit{arXiv:2508.14318}, 2025.

\bibitem{shoeybi2019megatron}
M.~Shoeybi \textit{et al.},
``Megatron-LM: Training multi-billion parameter language models using
model parallelism,''
\textit{arXiv:1909.08053}, 2019.

\bibitem{google2023tpuv4}
Google Cloud, ``TPU v4 enables performance, energy and CO$_2$e efficiency gains,''
2023.

\bibitem{epri2024powering}
J.~Aljbour \textit{et al.},
``Powering intelligence: Analyzing AI and data center energy consumption,''
EPRI White Paper 3002028905, 2024.

\end{thebibliography}

\end{document}
